\centerline{\bf \Huge Уравнения-Дроби-Модуль. Контрольная работа}
\centerline{\bf \Huge }

\begin{flushright}
        \scriptsize{\textbf{Для практики нет ничего более важного, чем хорошая теория. (Л. Больцман)}}
\end{flushright}

\problem{\textbf{1}} Решить уравнения с модулем:

а) $|X-15|=5$

б) $|111\dfrac{111}{754}\cdot X + 333\dfrac{112}{755}|=-123456789$

\textit{Знак \textbf{-}, это не всегда знак арифметического действия вычитания. Также он является одним целым с числом, когда число является отрицательным.}

\problem{\textbf{2}} Решить уравнение:

$2\dfrac{1}{4} : \dfrac{2}{3} = 1\dfrac{1}{2} : X$

\problem{\textbf{3}} При каком значении $A$ уравнение $A\cdot X = 36$:

а) Имеет корень, равный 3?

б) Имеет корень, равный 0,5?

в) Не имеет корней?